\subsection{TENSOR PRODUCT OF GRADED ALGEBRAS RELATIVE TO COMMUTATION FACTORS}

DEFINITION 6. Let \((\Delta_{i})_{i\in I}\) be a finite family of commutative monoids written additively. A system of commutation factors over the \(\Delta_{i}\) with values in a commutative ring A is a system of mappings \(\varepsilon_{ij}:\Delta_i\times \Delta_j\to \mathbf{A}\) , where \(i\in \mathbf{I}, j\in \mathbf{I}, i\neq j\) satisfying the following conditions:

\[
\varepsilon_ {i j} \left(\alpha_ {i} + \alpha_ {i} ^ {\prime}, \beta_ {j}\right) = \varepsilon_ {i j} \left(\alpha_ {i}, \beta_ {j}\right) \varepsilon_ {i j} \left(\alpha_ {i} ^ {\prime}, \beta_ {j}\right)\tag{22}
\]

\[
\varepsilon_ {i j} \left(\alpha_ {i}, \beta_ {j} + \beta_ {j} ^ {\prime}\right) = \varepsilon_ {i j} \left(\alpha_ {i}, \beta_ {j}\right) \varepsilon_ {i j} \left(\alpha_ {i}, \beta_ {j} ^ {\prime}\right)\tag{23}
\]

\[
\varepsilon_ {i j} (\alpha_ {i}, \beta_ {j}) \varepsilon_ {j i} (\beta_ {j}, \alpha_ {i}) = 1,\tag{24}
\]

for all \(\alpha_{i},\alpha_{i}^{\prime}\) in \(\Delta_{i},\beta_{j},\beta_{j}^{\prime}\) in \(\Delta_j\).

If I is given a total ordering and the \(\Delta_{i}\) are groups, a system of commutation factors is defined over the \(\Delta_{i}\) by taking for every ordered pair \((i,j)\) such that i\textless j an arbitrary Z-bilinear mapping of \(\Delta_{i}\times\Delta_{j}\) into the (multiplicative) Z-module A* of invertible elements of the ring A and then writing

\[
\varepsilon_ {j i} (\beta_ {j}, \alpha_ {i}) = (\varepsilon_ {i j} (\alpha_ {i}, \beta_ {j})) ^ {- 1}
\]

for \(i < j\) .

Note that, since the \(\varepsilon_{ij}(\alpha_{i},\beta_{j})\) are invertible,

\[
\varepsilon_ {i j} (0, \beta_ {j}) = \varepsilon_ {i j} (\alpha_ {i}, 0) = 1,
\]

by virtue of (22) and (23).

Examples. (1) The trivial system of commutation factors consists of the \(\varepsilon_{ij}\) such that \(\varepsilon_{ij}(\alpha_i, \beta_j) = 1\) for all \(i, j, \alpha_i \in \Delta_i, \beta_j \in \Delta_j\) .

\begin{enumerate}
\def\labelenumi{(\arabic{enumi})}
\setcounter{enumi}{1}
\tightlist
\item
  If we take \(\mathbf{A} = \mathbf{Z}\) and \(\Delta_{i} = \mathbf{Z}\) for all \(i \in \mathbf{I}\) , a system of commutation factors is obtained by taking \(\varepsilon_{ij}(\alpha_i, \beta_j) = (-1)^{\alpha_i\beta_j}\) . Note that this number depends only on the parities of \(\alpha_{i}\) and \(\beta_{j}\) and the \(\varepsilon_{ij}\) can therefore be considered as commutation factors when certain of the \(\Delta_{i}\) are equal to \(\mathbf{Z}/2\mathbf{Z}\) and the others to \(\mathbf{Z}\) .
\end{enumerate}

These two examples are the most frequent cases encountered in applications.

PROPOSITION 10. Let A be a commutative ring and \((\Delta_{i})_{i\in I}\) a finite family of commutative monoids written additively; for each \(i\in I\) let \(\mathbf{E}_i\) be a graded A-algebra of type \(\Delta_i\) . Finally, let \((\varepsilon_{ij})\) be a system of commutation factors over the \(\Delta_i\) with values in A. Then there exist a graded A-algebra E of type \(\Delta = \prod_{i\in I}\Delta_i\) and for each \(i\in I\) an algebra homomorphism \(h_i:\mathbf{E}_i\to \mathbf{E}\) , with the following properties:

\begin{enumerate}
\def\labelenumi{(\roman{enumi})}
\item
  If \(\phi_{i}:\Delta_{i}\to\Delta\) is the canonical homomorphism, then \(h_{i}\) is a graded homomorphism (II, § 11, no. 2), in other words, \(h_{i}(E_{i}^{\alpha_{i}})\subset E^{\phi_{i}(\alpha_{i})}\) , where \((E_{i}^{\alpha_{i}})\) and \((E^{\alpha})\) denote the respective graduations on \(E_{i}\) and \(E\) .
\item
  If \(i \neq j\) and \(x_{i}\) (resp. \(x_{j}\) ) is a homogeneous element of \(\mathbf{E}_i\) (resp. \(\mathbf{E}_j\) ) of degree \(\alpha_i \in \Delta_i\) (resp. \(\beta_j \in \Delta_j\) ), then
\end{enumerate}

\[
h _ {i} (x _ {i}) h _ {j} (x _ {j}) = \varepsilon_ {i j} (\alpha_ {i}, \beta_ {j}) h _ {j} (x _ {j}) h _ {i} (x _ {i}).\tag{25}
\]

\begin{enumerate}
\def\labelenumi{(\roman{enumi})}
\setcounter{enumi}{2}
\tightlist
\item
  For every A-algebra \(\mathbf{F}\) and every system of homomorphisms \(f_{i}:\mathbf{E}_{i}\to \mathbf{F}\) satisfying the conditions
\end{enumerate}

\[
f _ {i} (x _ {i}) f _ {j} (x _ {j}) = \varepsilon_ {i j} (\alpha_ {i}, \beta_ {j}) f _ {j} (x _ {j}) f _ {i} (x _ {i}),\tag{26}
\]

where \(i, j, x_i, x_j, \alpha_i, \beta_j\) are as in (ii), then there exists one and only one algebra homomorphism \(f: \mathbf{E} \to \mathbf{F}\) such that \(f_i = f \circ h_i\) for all \(i \in \mathbf{I}\) . Moreover the underlying A-module of \(\mathbf{E}\) is the tensor product \(\bigotimes_{i \in \mathbf{I}} \mathbf{E}_i\) .

Consider the A-module \(E = \bigotimes_{i \in I} E_i\) ; it is identified with the direct sum of the submodules \(E^\alpha\) , where, for each \(\alpha = (\alpha_i) \in \Delta\) , we write \(E^\alpha = \bigotimes_{i \in I} E_i^{\alpha_i}\) ; the \(E^\alpha\) therefore form a graduation of type \(\Delta\) on the A-module E. We shall define on E a graded A-algebra structure of type \(\Delta\) . For this let I be given a total ordering; for every ordered pair of elements \(\alpha = (\alpha_i)\) , \(\beta = (\beta_i)\) of \(\Delta\) , we must first define an A-bilinear mapping of \(E^\alpha \times E^\beta\) into \(E^\alpha + \beta\) , or alternatively an A-linear mapping \(m_{\alpha\beta}\) of \(E^\alpha \otimes_A E^\beta\) into \(E^\alpha + \beta\) . We shall define \(m_{\alpha\beta}\) by the condition

\[
m _ {\alpha \beta} \left(\left(\bigotimes_ {i \in I} x _ {i}\right) \otimes \left(\bigotimes_ {i \in I} y _ {i}\right)\right) = \varepsilon (\alpha , \beta) \bigotimes_ {i \in I} (x _ {i} y _ {i})\tag{27}
\]

for \(x_{i}\in \mathrm{E}_{i}^{\alpha_{i}},y_{i}\in \mathrm{E}_{i}^{\beta_{i}}\) , where

\[
\varepsilon (\alpha , \beta) = \prod_ {i > j} \varepsilon_ {i j} (\alpha_ {i}, \beta_ {j}).\tag{28}
\]

The right hand side of (27) obviously belongs to \(\mathbf{E}^{\alpha +\beta}\) and the mapping \((x_{1},\ldots ,x_{n},y_{1},\ldots ,y_{n})\mapsto \varepsilon (\alpha ,\beta)\bigotimes_{i\in I}(x_{i}y_{i})\) is A-multilinear in the product of the \(\mathrm{E}_i^{\alpha_i}\) and the \(\mathrm{E}_i^{\beta_i}(1\leqslant i\leqslant n)\) . Then it must be proved that the multiplication thus defined on E is associative; now, if \(\gamma = (\gamma_{i})\) is a third element of \(\Delta\) and \(z_{i} \in \mathrm{E}_{i}^{\gamma_{i}}\) for \(1 \leqslant i \leqslant n\) , then

\[
\left(\left(\bigotimes_ {i} x _ {i}\right) \left(\bigotimes_ {i} y _ {i}\right)\right) \left(\bigotimes_ {i} z _ {i}\right) = \varepsilon (\alpha + \beta , \gamma) \varepsilon (\alpha , \beta) \bigotimes_ {i} (x _ {i} y _ {i} z _ {i})
\]

\[
\left(\bigotimes_ {i} x _ {i}\right) \left(\left(\bigotimes_ {i} y _ {i}\right) \left(\bigotimes_ {i} z _ {i}\right)\right) = \varepsilon (\alpha , \beta + \gamma) \varepsilon (\beta , \gamma) \bigotimes_ {i} (x _ {i} y _ {i} z _ {i})
\]

and it reduces to verifying the identity

\[
\varepsilon (\alpha + \beta , \gamma) \varepsilon (\alpha , \beta) = \varepsilon (\alpha , \beta + \gamma) \varepsilon (\beta , \gamma).
\]

But the latter follows immediately from the relations

\[
\varepsilon (\alpha + \beta , \gamma) = \varepsilon (\alpha , \beta) \varepsilon (\beta , \gamma)
\]

\[
\varepsilon (\alpha , \beta + \gamma) = \varepsilon (\alpha , \beta) \varepsilon (\alpha , \gamma)
\]

themselves immediate consequences of the definition (28) and (22) and (23).

If, for all \(i \in \mathbf{I}\) , \(e_i\) denotes the unit element of \(\mathbf{E}_i\) , we know that \(e_i\) is homogeneous of degree 0 (§ 3, no. 1), hence \(e = \bigotimes_{i \in \mathbf{I}} e_i\) is homogeneous of degree 0 and it follows from (27), (28) and the relations

\[
\varepsilon_ {i j} (\alpha_ {i}, 0) = \varepsilon_ {i j} (0, \beta_ {j}) = 1
\]

that \(e\) is unit element of \(\mathbf{E}\) , which completes the definition on \(\mathbf{E}\) of the desired graded A-algebra structure. Then take \(h_i(x_i) = x_i \otimes \bigotimes_{j \neq i} e_j\) ; to verify that \(h_i(x_i x_i') = h_i(x_i) h_i(x_i')\) for \(x_i, x_i'\) in \(\mathbf{E}_i\) , attention may be confined to the case where \(x_i\) and \(x_i'\) are homogeneous and then this relation follows immediately from (27) and the relations \(\varepsilon_{ij}(\alpha_i, 0) = \varepsilon_{ij}(0, \beta_j) = 1\) ; the same relations and (24) prove also that the \(h_i\) satisfy conditions (i) and (ii) of the statement and that

\[
\bigotimes_ {i \in I} x _ {i} = \prod_ {i \in I} h _ {i} (x _ {i})\tag{29}
\]

where the right hand side is the product of the ordered sequence \((h_i(x_i))_{i\in I}\) in E with the given total ordering on I (I, § 1, no. 2) (it suffices to argue by induction on the number of \(x_{i}\) (assumed homogeneous) distinct from the \(e_i\) ).

It remains to prove condition (iii); note that the mapping

\[
(x _ {i}) _ {i \in I} \mapsto \prod_ {i \in I} f _ {i} (x _ {i}),
\]

where the right hand side is the product of the ordered sequence \((f_{i}(x_{i}))_{i\in\mathbf{I}}\) with the given total ordering on I, is A-multilinear. Then there exists one and only one A-linear mapping \(f: E \rightarrow F\) such that

\[
f \left(\bigotimes_ {i \in I} x _ {i}\right) = \prod_ {i \in I} f _ {i} (x _ {i}).\tag{30}
\]

Clearly \(f(e)\) is the unit element of \(\mathbf{F}\) and \(f \circ h_i = f_i\) ; to verify that \(f\) is an algebra homomorphism, in other words that \(f(x)f(y) = f(xy)\) for \(x, y\) in \(\mathbf{E}\) , attention may be confined, by linearity, to the case where \(x = \bigotimes_{i \in \mathbf{I}} x_i\) and \(y = \bigotimes_{i \in \mathbf{I}} y_i\) , \(x_i\) (resp. \(y_i\) ) being homogeneous of degree \(\alpha_i\) (resp. \(\beta_i\) ) for all \(i \in \mathbf{I}\) . The relation to be verified then reduces, by (27), to

\[
\left(\prod_ {i \in \mathrm{I}} f _ {i} (x _ {i})\right) \left(\prod_ {i \in \mathrm{I}} f _ {i} (y _ {i})\right) = \varepsilon (\alpha , \beta) \prod_ {i \in \mathrm{I}} \left(f _ {i} (x _ {i}) f _ {i} (y _ {i})\right).
\]

But, taking account of relations (26), this is a consequence of Lemma 2 of no. 6.

Clearly the algebra E and the canonical mapping \(\psi: \bigotimes_{i \in I} E_i \to E\) constitute a solution of the universal mapping problem (Set Theory, IV, § 3, no. 1), where \(\Sigma\) is the species of A-algebra structure and the \(\alpha\) -mappings \(\prod_{i} f_i\) from \(\prod_{i} E_i\) to an A-algebra, satisfying conditions (26).

For a fixed total ordering on I, the graded algebra E defined in the proof of Proposition 10 will be called a graded tensor \(\varepsilon\) -product of type \(\Delta\) of the family \((\mathrm{E}_i)_{i\in \mathbf{I}}\) of graded algebras of type \(\Delta_i\) and will be denoted by \({}^{\varepsilon}\bigotimes_{i\in \mathbf{I}}\mathrm{E}_i\) (if no confusion can arise over the ordering on I); similarly, the homomorphism \(f:\mathrm{E}\to \mathrm{F}\) defined in the proof of Proposition 10 will be denoted by \({}^{\varepsilon}\bigotimes_{i\in \mathbf{I}}f_i\) . The homomorphisms \(h_i\) are called canonical. We also write \({}^{\varepsilon}\mathrm{G}^{\otimes n}\) when \(\mathbf{I} = [1,n]\) and all the \(\mathrm{E}_i\) are equal to the same algebra G.

Remarks. (1) We recover the tensor product of algebras defined in no. 1 (with moreover the graduation the tensor product of those of its factors) taking \(\varepsilon_{ij}(\alpha_i, \beta_j) = 1\) for all \(i, j, \alpha_i\) and \(\beta_j\) .

\begin{enumerate}
\def\labelenumi{(\arabic{enumi})}
\setcounter{enumi}{1}
\tightlist
\item
  Suppose that all the \(\Delta_{i}\) are equal to \(\mathbf{Z}\) and write \(\varepsilon_{ij}(\alpha_i, \beta_j) = (-1)^{\alpha_i \beta_j}\) ; the tensor \(\varepsilon\) -product \({}^{\varepsilon}\bigotimes_{i \in I} E_i\) corresponding to this system of commutation factors is then called the skew tensor product of the graded algebras \(E_i\) of type \(\mathbf{Z}\) and denoted by \({}^{g}\bigotimes_{i \in I} E_i\) (or \(E{}^{g} \otimes_A F\) for two algebras, or \({}^{g} G^{\otimes n}\) instead of \({}^{\varepsilon} G^{\otimes n}\) ).
\end{enumerate}

COROLLARY 1. In the notation of Proposition 10, suppose further that F is a graded A-algebra of type \(\Delta\) and that each \(f_{i}\) is a graded algebra homomorphism relative to \(\phi_{i}:\Delta_{i}\to\Delta\) ; then \(f={}^{\varepsilon}\bigotimes_{i}f_{i}\) is a graded algebra homomorphism.

This follows immediately from the definition of \(f\) and the fact that

\[
\sum_ {i \in I} \phi_ {i} (\alpha_ {i}) = (\alpha_ {i})
\]

by definition of the \(\phi_{i}\) .

It is therefore seen that \((\mathrm{E},\psi)\) is also a solution of another universal mapping problem, where this time \(\Sigma\) is the species of graded A-algebra structure of type \(\Delta\) , the morphisms being graded algebra homomorphisms of type \(\Delta\) and the \(\alpha\) -mappings the mappings \(\prod_{i}f_{i}\) , where, in addition to conditions (26), it is assumed that \(f_{i}\) is a graded algebra homomorphism relative to \(\phi_{i}\) .

COROLLARY 2. Let \((\mathbf{E}_i)_{i\in \mathbf{I}}\) , \((\mathbf{F}_i)_{i\in \mathbf{I}}\) be two finite families of A-algebras, with \(\mathbf{E}_i\) and \(\mathbf{F}_i\) graded of type \(\Delta_i\) for all \(i\in \mathbf{I}\) . For each \(i\in \mathbf{I}\) , let \(g_i:\mathbf{E}_i\to \mathbf{F}_i\) be a graded algebra homomorphism of type \(\Delta_i\) . Then, if \(h_i:\mathbf{E}_i\to {}^{\varepsilon}\bigotimes_{i\in \mathbf{I}}\mathbf{E}_i\) and \(h_i':\mathbf{F}_i\to {}^{\varepsilon}\bigotimes_{i\in \mathbf{I}}\mathbf{F}_i\) are the canonical homomorphisms, there exists one and only one homomorphism of graded A-algebras of type \(\Delta\) , \(g: {}^{\varepsilon}\bigotimes_{i\in \mathbf{I}}\mathbf{E}_i\to {}^{\varepsilon}\bigotimes_{i\in \mathbf{I}}\mathbf{F}_i\) such that \(g\circ h_i = h_i'\circ g_i\) for all \(i\in \mathbf{I}\) . Also, if each \(g_i\) is bijective, so is \(g\) .

It suffices to apply Corollary 1 to \(f_{i} = h_{i}^{\prime} \circ g_{i}\) , noting that conditions (26) then follow from relations (25) applied to the \(h_{i}^{\prime}\) .

The homomorphism defined in Corollary 2 is also denoted by \({}^{\varepsilon}\bigotimes_{i} g_{i}\) (if no confusion can arise); if, for each \(i \in I\) , \(G_{i}\) is a third graded A-algebra of type \(\Delta_{i}\) and \(g_{i}^{\prime}: F_{i} \to G_{i}\) a graded algebra homomorphism, then

\[
\left(^ {\varepsilon} \bigotimes_ {i} g _ {i} ^ {\prime}\right) \circ \left(^ {\varepsilon} \bigotimes_ {i} g _ {i}\right) = ^ {\varepsilon} \bigotimes_ {i} \left(g _ {i} ^ {\prime} \circ g _ {i}\right),\tag{31}
\]

as follows immediately from (30).

In the case of a skew tensor product of graded algebras of type \(\mathbf{Z}\) , we write \({}^{g}\bigotimes_{i}f_{i}\) instead of \({}^{\varepsilon}\bigotimes_{i}f_{i}\) for homomorphisms \(f_{i}: \mathrm{E}_{i} \to \mathrm{F}_{i}\) of graded algebras of type \(\mathbf{Z}\) ; when \(\mathrm{I} = \{1, 2\}\) , this homomorphism is also denoted by \(f_{1}{}^{g}\otimes f_{2}\) ; when \(\mathrm{I} = [1, n]\) and all the \(\mathrm{E}_{i}\) (resp. \(\mathrm{F}_{i}\) ) are equal and all the \(f_{i}\) equal to the same homomorphism \(f\) , we write \({}^{g}f^{\otimes n}\) .

Remark. In the proof of Proposition 10, a total ordering on I was used to define an algebra structure on the tensor product \(\bigotimes_{i\in I}E_i\) of the A-modules

\(E_{i}\) . If the ordering on I is changed, another multiplicative structure arises on \(\bigotimes_{i\in I}E_{i}\) , but the new algebra thus obtained is canonically isomorphic to the above, since both are solutions of the same universal mapping problem. For example, when \(I=\{1,2\}\) , the canonical isomorphism of the algebra \(E_{1}{}^{\varepsilon}\otimes_{A}E_{2}\) onto the algebra \(E_{2}{}^{\varepsilon}\otimes_{A}E_{1}\) maps \(x_{1}\otimes x_{2}\) to \(\varepsilon_{2,1}(\alpha,\beta)x_{2}\otimes x_{1}\) , where \(x_{1}\) is homogeneous of degree \(\alpha\) and \(x_{2}\) homogeneous of degree \(\beta\) .

Let J be a subset of I and, for each \(i \in J\) , consider the canonical homomorphism \(h_{i}: E_{i} \to {}^{\varepsilon}\bigotimes_{i \in I} E_{i} = E\) . By virtue of relations (25) a canonical homomorphism \(h: E' = {}^{\varepsilon}\bigotimes_{i \in J} E_{i} \to E\) is derived canonically (by Proposition 10) from these homomorphisms, such that, for all \(i \in J\) , \(h_{i}' = h \circ h_{i}\) , \(h_{i}'\) being the canonical homomorphism \(E_{i} \to E'\) . Taking the total ordering on J induced by that chosen on I, we obtain

\[
h \left(\bigotimes_ {i \in J} x _ {i}\right) = \prod_ {i \in I} h _ {i} \left(x _ {i}\right) = \bigotimes_ {i \in I} x _ {i} ^ {\prime}
\]

where the middle term is the product of the ordered sequence \((h_i(x_i))_{i\in J}\) and where, in the right hand term, \(x_{i}^{\prime} = x_{i}\) for \(i\in J\) , \(x_{i}^{\prime} = e_{i}\) for \(i\notin J\) .

PROPOSITION 11. (``associativity'' of the tensor \(\varepsilon\) -product). In the notation of Proposition 10, let \((J_{\lambda})_{\lambda \in L}\) be a partition of I and write \(\Delta_{\lambda}' = \prod_{i \in J_{\lambda}} \Delta_i\) for all \(\lambda \in L\) . Let \(E_{\lambda}'\) be a graded tensor \(\varepsilon\) -product of type \(\Delta_{\lambda}'\) of the family \((E_i)_{i \in J_{\lambda}}\) (for some total ordering chosen on \(J_{\lambda}\) ). On the other hand, for \(\lambda, \mu\) in L and \(\lambda \neq \mu\) , we write, for \(\alpha_{\lambda}' = (\alpha_i)_{i \in J_{\lambda}}, \beta_{\mu}' = (\beta_j)_{j \in J_{\mu}}\) ,

\[
\varepsilon_ {\lambda \mu} ^ {\prime} (\alpha _ {\lambda} ^ {\prime}, \beta_ {\mu} ^ {\prime}) = \prod_ {i \in J _ {\lambda}, j \in J _ {\mu}} \varepsilon_ {i j} (\alpha_ {i}, \beta_ {j}).\tag{32}
\]

Then \((\varepsilon_{\lambda \mu}^{\prime})\) is a system of commutation factors over the \(\Delta_{\lambda}^{\prime}\) with values in A and there exists one and only one homomorphism of graded algebras of type \(\Delta, v: {}^{\varepsilon'}\bigotimes_{\lambda \in L} E_{\lambda}^{\prime} \to {}^{\varepsilon}\bigotimes_{i \in I} E_{i}\) , such that

\[
v \left(\bigotimes_ {\lambda \in L} \left(\bigotimes_ {i \in J _ {\lambda}} x _ {i}\right)\right) = \bigotimes_ {i \in I} x _ {i}\tag{33}
\]

for all \((x_{i})\in \prod_{i\in \mathbf{I}}\mathrm{E}_{i}\) , provided that the total ordering is taken on \(\mathbf{I}\) which induces on each \(\mathbf{J}_{\lambda}\) the chosen total ordering, and which is such that, for \(\lambda < \mu\) in \(\mathbf{L}\) , \(i\in \mathbf{J}_{\lambda}\) and \(j\in \mathbf{J}_{\mu}\) , \(i < j\) .

The fact that the \(\varepsilon_{\lambda \mu}'\) form a system of commutation factors is trivial. Let \(h_{i,\lambda}:\mathrm{E}_i\to \mathrm{E}_{\lambda}', h_{\lambda}':\mathrm{E}_{\lambda}'\to {}^{\varepsilon'}\bigotimes_{\lambda \in L}\mathrm{E}_{\lambda}'\) be the canonical homomorphisms (for \(\lambda \in L\) , \(i\in \mathbf{J}_{\lambda}\) ) and write \(h_i'' = h_\lambda' \circ h_{i,\lambda}\) ; it will suffice by virtue of the uniqueness of the solution of a universal mapping problem, to show that \({}^{\varepsilon'}\bigotimes_{\lambda \in L} E_{\lambda}'\) and the \(h_i''\) satisfy the conditions of Proposition 10. Now, for all \(\lambda \in L\) , let \(f_{\lambda}': E_{\lambda}' \to F\) be the unique algebra homomorphism such that \(f_{\lambda}' \circ h_{i,\lambda} = f_i\) for all \(i \in J_{\lambda}\) . We show that, for \(\lambda \neq \mu\) , \(\alpha_{\lambda}' = (\alpha_i)_{i \in J_{\lambda}}\) , \(\beta_{\mu}' = (\beta_j)_{j \in J_{\mu}}\) ,

\[
f _ {\lambda} ^ {\prime} (x _ {\lambda} ^ {\prime}) f _ {\mu} ^ {\prime} (x _ {\mu} ^ {\prime}) = \varepsilon_ {\lambda \mu} ^ {\prime} (\alpha_ {\lambda} ^ {\prime}, \beta_ {\mu} ^ {\prime}) f _ {\mu} ^ {\prime} (x _ {\mu} ^ {\prime}) f _ {\lambda} ^ {\prime} (x _ {\lambda} ^ {\prime})\tag{34}
\]

for \(x_{\lambda}^{\prime} \in \mathrm{E}_{\lambda}^{\prime}\) (resp. \(x_{\mu}^{\prime} \in \mathrm{E}_{\mu}^{\prime}\) ) homogeneous of degree \(\alpha_{\lambda}^{\prime}\) (resp. \(\beta_{\mu}^{\prime}\) ); it suffices, by linearity, to verify it when \(x_{\lambda}^{\prime} = \bigotimes_{i \in J_{\lambda}} x_{i}\) , \(x_{\mu}^{\prime} = \bigotimes_{j \in J_{\mu}} x_{j}\) , \(x_{i}\) (resp. \(x_{j}\) ) being homogeneous of degree \(\alpha_{i}\) (resp. \(\beta_{j}\) ) in \(E_{i}\) (resp. \(E_{j}\) ) for \(i \in J_{\lambda}, j \in J_{\mu}\) . But this follows from formula (30) which defines the \(f_{\lambda}^{\prime}\) and Lemma 3 of no. 6, taking account of hypothesis (26) and definition (32). There is therefore one and only one algebra homomorphism \(f: ^{\varepsilon'} \bigotimes_{\lambda \in L} E_{\lambda}' \to F\) such that \(f \circ h_{\lambda}' = f_{\lambda}'\) for all \(\lambda \in L\) ; whence \(f \circ h_{i}' = f_{i}\) for all \(i \in I\) and the uniqueness of \(f\) is trivial.

